\documentclass[11pt,a4paper]{report}

% ============================================================================
% RetailVision Edge AI — Documento de MVP (Producto Mínimo Viable)
% Universidad Santo Tomás, Seccional Tunja
% ============================================================================

\usepackage[utf8]{inputenc}
\usepackage[T1]{fontenc}
\usepackage{lmodern}
\usepackage[spanish,es-noindentfirst]{babel}
\usepackage[a4paper,margin=2.5cm]{geometry}
\usepackage{xcolor}
\usepackage{graphicx}
\usepackage{tabularx}
\usepackage{longtable}
\usepackage{booktabs}
\usepackage{enumitem}
\usepackage{fancyhdr}
\usepackage{titlesec}
\usepackage{hyperref}
\usepackage{array}
\usepackage{colortbl}
\usepackage{microtype}
\usepackage{amssymb}
\usepackage{textcomp}
\sloppy

% --- Paleta (coherente con el dashboard) ------------------------------------
\definecolor{rvAccent}{HTML}{C76A00}
\definecolor{rvAccentDark}{HTML}{8A4600}
\definecolor{rvInk}{HTML}{2C3242}
\definecolor{rvMuted}{HTML}{5B6478}
\definecolor{rvLine}{HTML}{CCD4E2}
\definecolor{rvCodeBg}{HTML}{F2F3F6}
\definecolor{rvGreen}{HTML}{1F9D66}
\definecolor{rvRed}{HTML}{D33B3B}

\titleformat{\chapter}[display]
  {\normalfont\bfseries\color{rvInk}}
  {\color{rvAccent}\Large\chaptertitlename\ \thechapter}
  {6pt}{\Huge}
\titleformat{\section}
  {\normalfont\Large\bfseries\color{rvInk}}{\thesection}{0.6em}{}
\titleformat{\subsection}
  {\normalfont\large\bfseries\color{rvAccentDark}}{\thesubsection}{0.6em}{}

\setlength{\headheight}{14.5pt}
\addtolength{\topmargin}{-0.5pt}
\pagestyle{fancy}
\fancyhf{}
\renewcommand{\headrulewidth}{0.4pt}
\renewcommand{\footrulewidth}{0pt}
\fancyhead[L]{\footnotesize\color{rvMuted} RetailVision Edge AI}
\fancyhead[R]{\footnotesize\color{rvMuted} \leftmark}
\renewcommand{\chaptermark}[1]{\markboth{#1}{}}
\fancyfoot[C]{\small\color{rvMuted} \thepage}

\hypersetup{
  colorlinks=true,
  linkcolor=rvAccentDark,
  urlcolor=rvAccentDark,
  citecolor=rvAccentDark,
  pdftitle={RetailVision Edge AI -- Documento de MVP},
  pdfauthor={Equipo RetailVision Edge AI -- Universidad Santo Tomas, Seccional Tunja},
}

\newcommand{\rvnota}[1]{%
  \par\noindent\fcolorbox{rvAccent}{rvCodeBg}{%
    \parbox{0.95\linewidth}{\small \textbf{\color{rvAccentDark}Nota:} #1}%
  }\par\vspace{4pt}
}

\newcommand{\rvcheck}{\textcolor{rvGreen}{\textbf{\checkmark}}}
\newcommand{\rvcross}{\textcolor{rvRed}{\textbf{\texttimes}}}

\setlist[itemize]{leftmargin=1.4em}
\setlist[enumerate]{leftmargin=1.6em}

\begin{document}

\begin{titlepage}
\centering
\vspace*{1.5cm}
{\Huge\bfseries\color{rvAccent} RetailVision Edge AI\par}
\vspace{0.4cm}
{\Large\color{rvInk} Documento de MVP (Producto Mínimo Viable)\par}
\vspace{0.3cm}
{\large\color{rvMuted} Sistema de Visión Computacional para \textit{Space Management}\\
y Analítica de Tráfico en Retail\par}
\vspace{2.2cm}

{\large\color{rvInk}\textbf{Universidad Santo Tomás, Seccional Tunja}\par}
\vspace{0.2cm}
{\normalsize\color{rvMuted} Desarrollo académico grupal\par}
\vspace{0.15cm}
{\normalsize\color{rvMuted} Docente asesor: Sergio Arley Puerto Moreno\par}

\vspace{2.2cm}
{\normalsize\color{rvInk}\textbf{Equipo}\par}
\vspace{0.3cm}
\begin{tabular}{l}
Silvana Sofía Siza Soriano \\
Yeikov Jared Ramos Sánchez \\
Sarita Toloza Buitrago \\
Juan David Sandoval Bello \\
Thomas Felipe Cepeda Mogollón \\
Wendy Mariana Vargas Ramírez \\
Andrés Felipe Galeano Téllez \\
Sara Valentina Caicedo Duarte \\
\end{tabular}

\vfill
{\large\color{rvMuted} 2026\par}
\end{titlepage}

\tableofcontents
\clearpage

% ============================================================================
\chapter{Resumen ejecutivo}
% ============================================================================

RetailVision Edge AI es el \textbf{Producto Mínimo Viable (MVP)} de un sistema de visión
computacional que convierte el video de las cámaras de seguridad de una tienda en decisiones de
\textit{space management}: detecta tomas y devoluciones de producto, mide el tráfico y la
permanencia por zona, y genera recomendaciones de redistribución de góndola, todo con
\textbf{privacidad por diseño} y procesamiento 100\,\% local.

Este documento describe el \textbf{problema que resuelve}, la \textbf{propuesta de valor}, el
\textbf{alcance exacto del MVP} (qué está construido y validado hoy, y qué queda fuera de esta
primera versión), la \textbf{evidencia de que funciona} (pruebas automatizadas y métricas de
rendimiento) y la \textbf{hoja de ruta} hacia una versión productiva.

\begin{center}
\renewcommand{\arraystretch}{1.3}
\begin{tabular}{>{\raggedright\arraybackslash}p{4.2cm} >{\raggedright\arraybackslash}p{9.8cm}}
\toprule
\rowcolor{rvCodeBg}\textbf{Estado del MVP} & \textbf{Funcional y validado} \\
\midrule
Pipeline de detección & \rvcheck\ Operativo (YOLO-pose + detector de producto + tracking) \\
Dashboard web & \rvcheck\ 5 vistas operativas + editor de zonas \\
API REST & \rvcheck\ 10 \textit{endpoints} documentados \\
Privacidad por diseño & \rvcheck\ Auditada (0 hallazgos en \texttt{qa\_privacidad.py}) \\
Pruebas automatizadas & \rvcheck\ 51 pruebas de backend + 5 de frontend, 100\,\% en verde \\
Despliegue en la nube & \rvcross\ Fuera de alcance de este MVP (ver hoja de ruta) \\
Multi-tienda / multi-tenant & \rvcross\ Fuera de alcance de este MVP \\
\bottomrule
\end{tabular}
\end{center}

% ============================================================================
\chapter{El problema}
% ============================================================================

\section{Contexto}

En un punto de venta físico, las decisiones de \textit{space management} (qué producto va en qué
góndola, con cuánto espacio) se toman hoy mayoritariamente \textbf{a ciegas}: con base en ventas
históricas y criterio del mercaderista, pero sin visibilidad de lo que ocurre \textbf{frente a la
góndola} en el momento de la decisión de compra.

\section{Lo que no se puede medir hoy, sin este sistema}

\begin{itemize}
  \item Cuántas personas se detienen frente a un estante y \textbf{no} compran.
  \item Cuántas \textbf{toman} un producto, lo consideran y lo \textbf{devuelven} (una señal de
    interés real que las ventas por sí solas no capturan).
  \item Qué zonas de la tienda concentran tráfico y cuáles pasan desapercibidas.
  \item El efecto de un cambio de planograma, sin esperar semanas de datos de venta.
\end{itemize}

\section{Por qué las soluciones existentes no bastan}

Las cámaras de conteo tradicionales miden tráfico agregado, no \textbf{interacción con el
producto}. Las soluciones de \textit{computer vision} de nivel empresarial que sí miden
interacción suelen requerir infraestructura en la nube, licenciamiento costoso, o procesar video
con reconocimiento facial --- lo que introduce un riesgo de cumplimiento normativo
(Ley 1581 de 2012 en Colombia) que muchas tiendas pequeñas y medianas no están dispuestas a asumir.

% ============================================================================
\chapter{Propuesta de valor}
% ============================================================================

\begin{center}
\renewcommand{\arraystretch}{1.25}
\begin{tabular}{>{\raggedright\arraybackslash}p{4.6cm} >{\raggedright\arraybackslash}p{9.4cm}}
\toprule
\rowcolor{rvCodeBg}\textbf{Para} & Tiendas y marcas de retail que necesitan optimizar el espacio
de góndola sin invertir en infraestructura de nube ni comprometer datos personales. \\
\textbf{Que necesitan} & Visibilidad objetiva de la interacción cliente-producto frente al
estante. \\
\textbf{RetailVision Edge AI} & Es un sistema de visión computacional \textit{edge} (local). \\
\textbf{Que} & Detecta tomas, devoluciones, tráfico y permanencia por zona, y recomienda
redistribuciones de planograma. \\
\textbf{A diferencia de} & Los conteos de tráfico genéricos o las soluciones de reconocimiento
facial en la nube. \\
\textbf{Nuestro sistema} & No identifica a nadie, procesa todo localmente y su motor de
recomendación es 100\,\% trazable (reglas sobre datos observados, no una caja negra). \\
\bottomrule
\end{tabular}
\end{center}

% ============================================================================
\chapter{Alcance del MVP}
% ============================================================================

\section{Funcionalidades incluidas (construidas y validadas)}

\begin{enumerate}
  \item \textbf{Ingesta de video}: subida desde el navegador, validación de formato y tamaño,
    extracción de vista previa anonimizada.
  \item \textbf{Calibración de zonas}: editor gráfico de polígonos (cuerpo + área de agarro) por
    góndola, con validación geométrica en vivo.
  \item \textbf{Pipeline de detección}: pose + tracking de personas (YOLO-pose + ByteTrack) y
    detección de producto en mano (modelo \textit{fine-tuneado}).
  \item \textbf{Motor de comportamiento}: máquina de estados que clasifica a cada persona
    (tránsito, llegada, inspección, mano en estante, toma, devolución).
  \item \textbf{Analítica de zonas}: heatmap de tráfico y clasificación caliente / templada /
    fría por góndola.
  \item \textbf{Agente de recomendación determinista}: texto accionable por zona y propuesta de
    reparto de \textit{facing} entre la zona más caliente y la más fría.
  \item \textbf{Dashboard web}: 5 vistas (Lectura general, Analítica, Plano de tienda,
    Evidencias, Informe) y un informe imprimible.
  \item \textbf{Privacidad por diseño}: anonimización de rostros dentro del propio pipeline,
    auditada automáticamente.
  \item \textbf{Accesibilidad}: panel de preferencias (contraste, tamaño de texto, lector de
    pantalla, navegación por teclado).
\end{enumerate}

\section{Explícitamente fuera de alcance de este MVP}

\begin{itemize}
  \item \textbf{Despliegue multi-tienda en la nube} con autenticación y aislamiento por cliente
    (\textit{multi-tenancy}). El MVP corre como instalación local de un solo punto de venta.
  \item \textbf{Integración directa con sistemas de punto de venta (POS)} o inventario para
    cruzar interacción con ventas reales.
  \item \textbf{Procesamiento de cámaras en vivo (streaming)}: el MVP procesa video ya grabado,
    no una transmisión continua.
  \item \textbf{Modelos de detección específicos por categoría de producto} más allá del dataset
    de referencia usado para el \textit{fine-tuning} actual.
  \item \textbf{Panel de administración de usuarios y roles} (el MVP no tiene autenticación:
    está pensado para operación interna de un equipo, no para exposición pública).
\end{itemize}

\rvnota{Estas exclusiones son decisiones deliberadas de alcance, no limitaciones técnicas
irresolubles: la arquitectura por capas del sistema (ver Manual Técnico, capítulo de
Arquitectura) fue diseñada precisamente para permitir añadir estas capacidades --- por ejemplo,
un repositorio de zonas en base de datos en vez de JSON, o un detector de personas distinto ---
sin reescribir el dominio.}

% ============================================================================
\chapter{Validación: evidencia de que el MVP funciona}
% ============================================================================

\section{Pruebas automatizadas}

\begin{center}
\begin{tabular}{>{\raggedright\arraybackslash}p{5.4cm} >{\raggedright\arraybackslash}p{3cm} >{\raggedright\arraybackslash}p{4.6cm}}
\toprule
\rowcolor{rvCodeBg}\textbf{Suite} & \textbf{Resultado} & \textbf{Qué cubre} \\
\midrule
Backend (\texttt{pytest}) & \rvcheck\ 51 / 51 & Dominio, orquestación, API \\
Frontend (\texttt{node:test}) & \rvcheck\ 5 / 5 & Lógica de tasas, zonas, Sankey \\
Auditoría de privacidad & \rvcheck\ 0 hallazgos & Ausencia de datos biométricos en salidas \\
Invariantes de métricas & \rvcheck\ OK & Embudo, heatmap, recomendaciones coherentes \\
\bottomrule
\end{tabular}
\end{center}

\section{Rendimiento}

Medido sobre GPU NVIDIA RTX 2060, video 720p: \textbf{\textasciitilde 24 FPS} de procesamiento,
es decir, un margen de \textbf{\textasciitilde 4$\times$} respecto al tiempo real (un video de un
minuto se procesa en \textasciitilde 15 segundos). Sin GPU, el sistema sigue siendo funcional en
CPU, a un FPS menor, sin cambios de código.

\section{Exactitud del embudo de interacción}

La detección de tomas y devoluciones se valida contra un \textit{ground truth} anotado
manualmente sobre el video de referencia, con precisión, \textit{recall} y F1 calculados por
\texttt{app/interfaces/validar\_embudo.py}, incluyendo un chequeo de coherencia automático
(ningún estante puede registrar tomas sin visitas, la tasa de rechazo no puede superar el
100\,\%, etc.).

% ============================================================================
\chapter{Cumplimiento normativo}
% ============================================================================

El MVP fue diseñado desde el primer día bajo el principio de \textbf{privacidad por diseño}, en
cumplimiento de la \textbf{Ley 1581 de 2012} (Colombia): anonimización de rostros antes de
cualquier render o persistencia, ausencia de reconocimiento facial y de identificadores
persistentes, y procesamiento 100\,\% local. El detalle técnico de esta garantía se documenta en
el Manual Técnico, capítulo de Privacidad.

% ============================================================================
\chapter{Limitaciones conocidas}
% ============================================================================

\begin{itemize}
  \item La precisión de la detección de producto-en-mano depende del dataset de
    \textit{fine-tuning} usado; productos muy distintos a los del dataset de referencia pueden
    requerir reentrenamiento.
  \item El sistema procesa \textbf{un análisis a la vez} por instancia (no hay cola de trabajos
    concurrentes en esta versión).
  \item Las zonas se calibran manualmente por video/cámara; un cambio de ángulo de cámara exige
    recalibrar.
  \item No hay autenticación de usuarios: el acceso al panel se asume restringido a la red
    interna del equipo que lo opera.
\end{itemize}

% ============================================================================
\chapter{Hoja de ruta}
% ============================================================================

\begin{center}
\renewcommand{\arraystretch}{1.3}
\begin{tabular}{>{\raggedright\arraybackslash}p{3.4cm} >{\raggedright\arraybackslash}p{10cm}}
\toprule
\rowcolor{rvCodeBg}\textbf{Horizonte} & \textbf{Próximos pasos propuestos} \\
\midrule
Corto plazo & Despliegue en un entorno de nube controlado para piloto con un cliente real;
ampliación del \textit{ground truth} de validación con más escenarios de tienda. \\
Mediano plazo & Autenticación básica y soporte multi-cámara por tienda; cola de análisis
concurrentes. \\
Largo plazo & \textit{Multi-tenancy} para operar varias tiendas/cadenas desde una sola
instalación; integración con sistemas de punto de venta para cruzar interacción con ventas
reales. \\
\bottomrule
\end{tabular}
\end{center}

% ============================================================================
\chapter{Conclusión}
% ============================================================================

RetailVision Edge AI demuestra, con un pipeline funcional, validado por pruebas automatizadas y
una auditoría de privacidad, que es posible medir la interacción real entre clientes y productos
en una góndola \textbf{sin} comprometer datos personales ni depender de infraestructura en la
nube. Este MVP establece la base arquitectónica (capas limpias, puertos y adaptadores) sobre la
que pueden construirse, sin reescritura, las capacidades de la hoja de ruta hacia una versión
productiva y multi-tienda.

\vspace{1cm}
\textbf{Equipo:} Silvana Sofía Siza Soriano, Yeikov Jared Ramos Sánchez, Sarita Toloza Buitrago,
Juan David Sandoval Bello, Thomas Felipe Cepeda Mogollón, Wendy Mariana Vargas Ramírez, Andrés
Felipe Galeano Téllez, Sara Valentina Caicedo Duarte.

\textbf{Docente asesor:} Sergio Arley Puerto Moreno --- Universidad Santo Tomás, Seccional
Tunja.

\end{document}
